\subsection{All rank-selected levels as exact kernel quotients}

\begin{result}{D.34}{Proved}{Every level measures an explicit quotient of the congruence kernel}
Fix a stable type $(g,n)$, put $d=3g-3+n$, and retain
$K=K_{g,n}$ from (D.40).  For an actual multicurve $\sigma$,
let $S_v$ be its cut components, with types $(g_v,n_v)$
including the individually pinned node branches.
There is a natural continuous homomorphism
\[
 \kappa_\sigma:\prod_v K_{g_v,n_v}
                  \longrightarrow K\subset\widehat\Gamma_{g,n}
                                                               \tag{D.55}
\]
defined by the cutting diagrams below.  Its image is exactly
$K\cap\overline{C_\sigma}$, where $C_\sigma$ is the actual
discrete multicurve stabilizer and the bar is closure in
$\widehat\Gamma_{g,n}$.

For $k\geq2$, let $J_k$ be the closed normal subgroup of $K$
generated by all the subgroups
\[
 h\,\operatorname{im}(\kappa_\sigma)h^{-1},
 \qquad h\in\widehat\Gamma_{g,n},\qquad
                         |\sigma|\geq d-k.                  \tag{D.56}
\]
Only one representative of each actual multicurve orbit is
needed in this formula.  The empty multicurve is included
when $k\geq d$; its map (D.55) is the identity of $K$.
Then there is a continuous group isomorphism, with the usual
inner ambiguity from the basepoint,
\[
               \Pi_1(\check P_k)\ \cong\ K/J_k.             \tag{D.57}
\]
Every finite continuous representation of the right side
corresponds explicitly to a finite covering of $\check P_k$
by the quotient construction in the proof.

For $2\leq k\leq6$ and $k<d$, all component kernels in
(D.55) vanish and (D.57) is exactly D.32.  For $k\geq d$
it is the trivial group.  At every intermediate level,
(D.55)--(D.57) specify the remaining kernel equation without
asserting its vanishing.  Inclusions of levels induce the
actual quotient maps $K/J_k\to K/J_{k+1}$, since
$J_k\subseteq J_{k+1}$.
\end{result}
\begin{proof}
We first construct (D.55), including the topology and its
image.  Write $C_{\vec\sigma}$ for the subgroup fixing
every component, branch and oriented curve.  The two
cutting rows are (D.46), now with arbitrary component
genera.  The top row is exact because the node-twist
lattice has its full induced topology by E.29 and
completion is right exact.  The bottom row is the
actual procongruence cutting sequence of the audited
Theorem 4.5 in \cite{ET_Boggi2014}.  The left vertical
map is the identity of the full embedded lattice
$\widehat{\mathbb Z}^{\sigma}$.  Hence restriction of
the upper cutting map gives an isomorphism
\[
 \ker\bigl(\widehat C_{\vec\sigma}
                         \longrightarrow\check C_{\vec\sigma}\bigr)
          \xrightarrow{\sim}\prod_v K_{g_v,n_v}.             \tag{D.58}
\]
Here is its inverse explicitly.  Lift a tuple of component
kernel elements to $c\in\widehat C_{\vec\sigma}$.  Its
image in $\check C_{\vec\sigma}$ belongs to the node-twist
lattice.  There is a unique $t$ in the identical upper
lattice having that image.  Send the tuple to $t^{-1}c$.
A different lift changes $c$ by the same twist as $t$,
so the result is independent of the lift.  The cutting
map restricted to the kernel is a continuous bijective
homomorphism of compact Hausdorff groups, proving both
the homomorphism law and continuity of this inverse.

Compose that inverse with the actual homomorphism
$\widehat C_{\vec\sigma}\to\widehat\Gamma_{g,n}$.
Its image lies in $K$ by commutativity with the lower
row, and this defines $\kappa_\sigma$.  To prove the
image assertion, let $x\in K\cap\overline{C_\sigma}$.
The compact image of $\widehat C_\sigma$ in the ambient
full group is precisely $\overline{C_\sigma}$; thus $x$
has a lift $c\in\widehat C_\sigma$.  Its image in
$\check C_\sigma$ is trivial because that group is an
actual subgroup of $\check\Gamma_{g,n}$.  The full and
procongruence permutation quotients have the identical
finite image $A_\sigma$ from (D.50).  Consequently $c$
belongs to the pinned subgroup and is in the kernel
of (D.58).  Thus $x$ is in the image of (D.55).
The reverse inclusion follows directly from its
construction.  This argument specifies the map even
when the full homomorphism
$\widehat C_\sigma\to\widehat\Gamma_{g,n}$ has a kernel.

Put $Y=P_k^{\mathrm{lev}}$ and $X=\check P_k$.
There are finitely many orbits of vertices and of
simplices.  On each such orbit, for the actual flag
stabilizer $C$, the two completions are
\[
        \widehat\Gamma/\overline C,
              \qquad \check\Gamma/\check C.
\]
The quotient of the first by $K$ is the second:
the image of the compact group $\overline C$ is
exactly $\check C$.  All face and degeneracy maps
are induced by the same discrete flag maps.  We
therefore have a degreewise quotient of simplicial
profinite sets
\[
                         X=K\backslash Y.                  \tag{D.59}
\]
Its vertex stabilizers in $K$ are exactly the groups
in (D.56), by (D.55) and translation of each orbit.
D.11 proves that $Y$ is connected and has trivial
profinite fundamental group.  We now compute finite
coverings of (D.59) directly; freeness of its vertex
action is not required.

Let $E\to X$ be a finite profinite covering.  Pull it
back to $Y$.  The covering classification of
\cite[\S2.1, Proposition 2.2]{D_Quick2008} and
$\Pi_1(Y)=1$ give an isomorphism
$Y\times_X E\simeq Y\times S$, where $S$ is a finite
set.  The canonical $K$-action on the pullback fixes
the $E$-coordinate and acts on the $Y$-coordinate.
In the chosen trivialization it has the form
\[
       a\cdot(y,s)=(ay,\alpha(a,y)s),
       \qquad \alpha(a,y)\in\operatorname{Sym}(S).           \tag{D.60}
\]
For fixed $a$, this permutation is a continuous
simplicial function of $y$ into the constant finite
set $\operatorname{Sym}(S)$.  Such a function factors
through $\pi_0(Y)$; connectedness makes it constant.
Write it as $\alpha(a)$.  Evaluation at one vertex
shows that $a\mapsto\alpha(a)$ is continuous, and
the action law in (D.60) makes it a homomorphism.
When $a$ fixes a vertex $y$, it fixes every point
$(y,e)$ of the canonical pullback.  Its action on
$S$ is therefore trivial.  Every vertex stabilizer,
and hence $J_k$, is in the kernel of $\alpha$.
Thus a finite cover gives a finite continuous
$K/J_k$-set.

Conversely, let $S$ be a finite continuous $K$-set
on which $J_k$ acts trivially.  Form the degreewise
compact quotient
\[
                       E_S=(Y\times S)/K\longrightarrow X.
\]
These are profinite spaces.  For example the orbit
relation of a continuous compact group action on
a profinite set is closed, and invariant clopen
neighbourhoods separate its orbits: for an invariant
open neighbourhood of a compact orbit, choose a
clopen neighbourhood of that orbit inside it and
take the complement of the saturation of its
complement.  That saturation is compact and is
open as a union of translates of a clopen set.
This gives an invariant clopen
neighbourhood inside the original one.  Hence the
quotient is compact Hausdorff and totally disconnected.

We check the precise covering property.  Choose
a simplex $\eta$ of $X$ and a specified lift of
one of its vertices to $E_S$.  Represent that
vertex by $(y,s)$.  A lift of $\eta$ to $Y$ can
be translated by an element of $K$ so that its
chosen vertex is $y$, because that vertex fibre
is one orbit.  Pair this lift with $s$ and take
its class in $E_S$.  A second such choice differs
by a member of the stabilizer of $y$; that member
acts trivially on $s$.  The resulting simplex of
$E_S$ is therefore independent of the choice
and is the unique possible lift.  The fibre at
each vertex is $S$, since its stabilizer acts
trivially there.  The degreewise comparison with
the pullback determined by a vertex is a
continuous bijection between compact Hausdorff
spaces, hence a homeomorphism.  This is exactly
the finite-covering definition in Quick's
$\Delta[0]\to\Delta[m]$ lifting criterion.

The constructions are inverse on covers and
on their maps.  Explicitly the canonical
quotient $(Y\times_X E)/K\to E$ is a compact
bijection in each degree: $K$ acts transitively
on every fibre of $Y\to X$ and fixes the
$E$-coordinate.  In the other direction the
map $Y\times S\to Y\times_X E_S$ sending
$(y,s)$ to $(y,[y,s])$ is a degreewise
compact bijection; its injectivity is exactly
the trivial action of vertex stabilizers,
and the same statement in each simplex
follows from the unique lift just proved.
Covering maps correspond under these maps
to the equivariant maps of finite sets.

We have obtained an equivalence of finite
covering categories which identifies their
fibre functors after choosing a vertex of
$Y$ above the base vertex of $X$.  Their
automorphism groups are therefore exactly
(D.57), as topological profinite groups.
Changing the lift translates the fibre by
an element of $K$ and gives the stated inner
ambiguity.  This proves the formula using
finite covers, without a higher-homotopy
inverse-limit assertion.

For $2\leq k\leq6$ and $k<d$, the component
dimension calculation (D.45) puts every
component in genus at most two.  D.19 then
makes every factor in (D.55) trivial.  When
$k\geq d$, the empty multicurve gives the
identity map on $K$, and $J_k=K$.  Finally
increasing $k$ only adds vertex orbits to
(D.56).  Pullback of coverings along the
inclusion of levels is restriction of the
same finite $K$-set, so its fundamental-group
map is precisely $K/J_k\to K/J_{k+1}$.
This proves all the claimed comparisons.
\end{proof}
